% ============================================================
% chapters/08-summary.tex -- פרק 8: סיכום – מוכנים לשלב הביצוע
% Content only. All design comes from template.cls.
% Figure: figures/ch8-next-steps.tex (TikZ)
% ============================================================

\chapter{סיכום: מוכנים לשלב הביצוע}
\label{ch:summary}

\paragraph{ממצאי ההכנה.} הדוח מציע לבחון את הקשר בין התנצרות, השתלבות הכנסייה ויציבות מדינית באמצעות השוואה מדידה של המערב הרומי, המזרח והממלכה הפרנקית. בספרות שנכללה בסקירה לא זוהתה השוואה התואמת בדיוק את המדדים המוצעים; זהו פער במסגרת הסקירה, ולא הוכחה שאין מחקר קודם בנושא (פרקים~\LR{\ref{ch:problem}--\ref{ch:literature}}). הכלים שנסקרו מספקים חלק מן התשתית, אך זמינות המקורות והיתר השימוש בכל מהדורה ייבדקו לפני איסוף הקורפוס (פרק~\ref{ch:market}). מטריצת ההחלטה מעדיפה שילוב של מאגר נתונים וניתוח טקסט בעזרת \en{AI}, עם מבחן מומחה חתום ונפרד מקטעי הפיתוח. בדיקת הרגישות תומכת ביציבות ההעדפה תחת ההנחות שנבדקו, ואינה מדד להצלחת המחקר בעתיד (פרק~\ref{ch:alternatives}).

\paragraph{המוכנות לשלב הבא.} שבע הדרישות כוללות ערכי יעד ושיטות אימות. התוכנית כוללת חמש חבילות עבודה, שלוש אבני דרך ותקציב תכנון לצורך ההצעה: כ־111 אלף ש״ח כלכליים, וכ־22 אלף ש״ח במזומן אם זמן החוקרים ניתן כתרומה אקדמית. זהו אומדן ולא התחייבות למימון. המעבר לביצוע מותנה בזמינות מומחים ומקורות ובאישור התוכנית (פרק~\ref{ch:workplan}). הפחתת ארבעת הסיכונים הקריטיים היא תחזית תכנון המותנית בפעולות המנע; הכיסוי, הדיוק והעלויות טרם נמדדו (פרק~\ref{ch:risks}).

\paragraph{שלוש הפעולות המיידיות.} לאחר אישור הדוח יבוצעו הפעולות שבאיור~\ref{fig:next-steps}: גיוס מומחה ומומחה גיבוי וקביעת מועד התיקוף; הכנת ספר קידוד ורשימת מקורות; ומיני־פיילוט נפרד על 50 קטעים לבדיקת זמינות, עלות ואיכות ראשונית. מיני־הפיילוט אינו מחליף את מבחן 300 הקטעים החתום. המחקר עצמו לא בוצע במסגרת מטלה זו.

\begin{figure}[htbp]
  \centering
  \input{figures/ch8-next-steps}
  \caption{שלוש הפעולות המיידיות לאחר אישור הדוח}
  \label{fig:next-steps}
\end{figure}

\begin{keybox}{בחזרה לאוגוסט 410}
כשהוויזיגותים נכנסו לרומא, בני התקופה מיהרו לשאול מי אשם: האלים שננטשו, או האל החדש. המחקר המוצע לא מבקש להכריע מי צדק אז, אלא לאפשר לבדוק את השאלה באופן שיטתי – בנתונים, בשיטה שאפשר לשחזר, ובאותם מדדים לשלוש מדינות שבחרו באותה דת והגיעו לגורלות שונים כל כך.
\end{keybox}
